\scanpage{15}
\deleted{\uline{type} $\underline{R}$ \uline{de} $S$ une section cartésienne de $\underline{D}$ sur $\underline{C}$. Ces revêtements de type $\underline{R}$ forment une catégorie \inserted{que nous noterons $\mathrm{Rev}_{\underline{R}}(S)$.} \handwritten{(} notée suivant les notations standard $\underset{\handwritten{\longleftarrow}}{\mathrm{Lim}}\,\underline{D}/\underline{C}$\handwritten{)}. On définira la \uline{réalisation géométrique} d'un revêtement $\handwritten{\underline{\phi}}$ de type $\underline{R}$ comme étant la limite inductive (si elle existe) dans $(\mathrm{Sch}$ du foncteur composé $\underline{C}\overset{\handwritten{\underline{\phi}}}{\qquad}\underline{D}\qquad(\mathrm{Sch})_{/S}$, où le deuxième de ces foncteurs est le foncteur source \deleted{\handwritten{(}associant à tout \unsure{revêtement}}(dé\inserted{d}uit du foncteur source sur la catégorie des schémas finis étales relatifs). Si pour tout $\handwritten{\underline{\phi}}$, cette limite inductive existe,} on obtient un foncteur réalisation géométrique
\[
\deleted{\handwritten{\unsure{\phi|r}}}\quad r:\mathrm{Rev}_{\underline{R}}(S)\qquad(\mathrm{Sch})_{/S}\,.
\]
Dans les cas qui nous intéresseront par la suite, ce foncteur $r$ sera fidèle, mais pas nécessairement pleinement fidèle; il importera donc de bien distinguer entre une revêtement de type $\underline{R}$ de $S$, et sa réalisation géométrique, le premier n'étant pas déterminée par la seule connaissance du second.

Avant d'aller plus loin, donnons quelques exemples.

\uline{Exemple} 1. Donnons-nous un schéma relatif $Z$ sur $S$, muni d'un groupe fini $G$ opérant \inserted{(à droite disons)} par $S$-automorphismes. Soit $C_{\unsure{o|0}}$ la catégorie réduite à un objet $\inserted{e}$ et dont le monoïde des endomorphismes est $G$, et prenons pour $\underline{C}$ la catégorie produit $\underline{\mathrm{Et}}_S\mathrm{x}C_o$, \deleted{pour foncteur $\underline{s}$ le foncteur $(S',e)$ qui aux les objets} La donnée d'un foncteur cartésien de \deleted{Et} $\underline{C}$ dans une catégorie fibrée $\underline{C}'$ sur $\underline{\mathrm{Et}}_S$ équivaut manifestement à la donnée d'un objet de la catégorie fibre $\underline{C}'_S$ avec une opération de $G$ dessus, de sorte que $(Z,G)$ nous \unsure{fournit} bien un foncteur cartésien $\underline{s}$. On trouve aussitôt que la catégorie $\mathrm{Rev}_{\underline{R}}(S)$ est équivalente canoniquement à la \deleted{\unsure{cat}} catégorie des revêtements finis ét\unsure{q|a}les $Z'$ de $Z$, munis d'une opération (droite de $G$ \deleted{compatible} rendant le morphisme structural $Z'\quad Z$ un $G$-morphisme.
% archive sheet number: 14
